\documentclass[12pt,a4paper]{article}

% =================================================
% Packages
% =================================================

\usepackage{amsmath,amssymb,amsfonts}
\usepackage{graphicx}
\usepackage{tikz}
\usepackage{subfig}
\usepackage{algpseudocode}
\usepackage{algorithm}
\usepackage{booktabs}
\usepackage{tabularx}
\usepackage{textcomp}
\usepackage{xcolor}
\usepackage{fontspec}
\usepackage{listings}
\usepackage{float}
\usepackage{placeins}
\usepackage[normalem]{ulem}
\usepackage{lscape}
\usepackage{geometry}
\usepackage{setspace}

% =================================================
% Page Setup
% =================================================

\geometry{
    a4paper,
    left=2.5cm,
    right=2.5cm,
    top=2.5cm,
    bottom=2.5cm
}

% =================================================
% Thai Font
% =================================================

\XeTeXlinebreaklocale "th"
\XeTeXlinebreakskip = 0pt plus 1pt

\setmainfont[
    BoldFont={THSarabunNew_Bold.ttf},
    ItalicFont={THSarabunNew_Italic.ttf},
    BoldItalicFont={THSarabunNew_BoldItalic.ttf}
]{THSarabunNew.ttf}

% =================================================
% Line Spacing
% =================================================

\setstretch{1.15}

% =================================================
% Section Formatting
% =================================================

\usepackage{titlesec}

% Section
\titleformat{\section}
    {\centering\Large\bfseries}
    {Section \thesection}
    {0.5em}
    {}

% Subsection
\titleformat{\subsection}
    {\large\bfseries}
    {\thesubsection}
    {0.5em}
    {}

% =================================================
% Document
% =================================================

\begin{document}


% =================================================
% SECTION 1
% =================================================

\section{ }
\subsection{Project introduction \& problem statement} 

In recent years, unmanned aerial vehicles (UAVs) have
become increasingly important in search and rescue,
disaster response, and emergency operations. UAVs can
provide aerial information over large areas and access
locations that may be difficult or dangerous for humans
to reach.

In emergency situations, locating humans quickly is an
important task. Traditional search methods may require
significant amounts of time and personnel, particularly
when the search area is large or difficult to access.
Therefore, an automated system capable of detecting and
localizing possible humans from UAV data can help improve
the efficiency of search operations.

This project proposes an AI-based human detection and
localization system capable of identifying possible human
locations using UAV camera and sensor data. The system
processes information collected by the UAV and estimates
the location of detected humans.

The detected information can then be reported to a
monitoring system or ground station, allowing operators
to view and analyze possible human locations.


% =================================================
% 1.1 Objectives
% =================================================

\subsection{Objectives}



\begin{enumerate}

    \item To develop an AI-based system capable of
    detecting humans from UAV camera data.

    \item To estimate and report the location of detected
    humans using UAV sensor information.

    \item To develop a system capable of processing
    detection information efficiently.

    \item To provide a monitoring interface for displaying
    UAV and human detection information.

    \item To evaluate the performance of the proposed
    human detection and localization system.

\end{enumerate}


% =================================================
% 1.2 Scope
% =================================================

\subsection{Scope of the Project}

The scope of this project covers the development of a
UAV-based human detection and localization system. The
system consists of UAV hardware, camera and sensor
components, an AI-based detection module, and a monitoring
system.

The system focuses on detecting possible humans within
the UAV camera field of view and estimating their
locations based on available sensor and positioning data.


% =================================================
% 1.3 Expected Results
% =================================================

\subsection{Expected Results}

The expected result of this project is a prototype system
that can detect possible humans from UAV camera data and
provide their estimated locations to an operator.

The system is expected to support search and rescue
operations by reducing the time required to identify
possible human locations and providing useful information
for further investigation.


% =================================================
% SECTION 2
% =================================================

\clearpage


% =================================================
% END
% =================================================

\end{document}